\section{H4: What the settlement price decides}\label{sec:derivatives}

Whether anyone has a reason to move the settlement price depends on how the derivative is
settled. This section examines the stock derivatives directly, on the \ConvSessions{} sessions
for which the derivatives trade files are available, and returns to the index at the end.

\subsection{The expiring future is priced on the shares}

Under settlement by delivery, the settlement price of an expiring stock future divides the
cash flows of a long position between the final variation margin and the payment for the
delivered shares, but does not change their sum. A long position bought at price $F$ pays
$F$ for the shares whatever the settlement price turns out to be. The expiring future should
therefore be priced on the value of the shares, that is on the current cash price, until the
end of the window, rather than on the average that determines the settlement price. If it were
settled in cash against the average, it would instead converge to the average, which by the
last minutes of the window is almost entirely determined.

Table~\ref{tab:convergence} distinguishes the two. The futures price and the cash price are
both measured relative to the settlement price, and the first is regressed on the second. By
the minute beginning two minutes before the close, \ConvExpiryMTwoEightFixed{}\% of the window's volume has already
traded, so the settlement price is nearly fixed. A future priced on the settlement price would
not respond to the cash price at that point; a future priced on the shares would move with it
one for one. The slope is \ConvExpiryMTwoEight{}, and over the final five minutes
\ConvExpiryLastFive{}, both far from zero and close to one. The median distance of the futures
price from the settlement price in that minute, \ConvExpiryMTwoEightGapsettle{} basis points, exceeds
its distance from the concurrent cash price, \ConvExpiryMTwoEightGapcash{} basis points.
Figure~\ref{fig:convergence} shows the same comparison minute by minute.

\begin{table}[htbp]
\centering
\small
\caption{Pricing of the expiring stock future in the settlement window. Regression of the
futures price on the cash price, both in basis points relative to the settlement price, over
security-minutes of the window; a slope of one indicates that the future moves with the cash
price, a slope of zero that it is priced on the settlement price. Standard errors clustered by
security in parentheses. Control sessions are the trading day before each expiry, on which the
same future has one day to run.}
\label{tab:convergence}
\input{generated/tables/convergence}
\end{table}

\begin{figure}[htbp]
  \centering
  \includegraphics[width=0.9\textwidth]{convergence.pdf}
  \caption{Median absolute distance of the expiring future's price from the concurrent cash
    price and from the settlement price, by minute of the window.}
  \label{fig:convergence}
\end{figure}

The result has a direct implication for the settlement incentive. A holder of an expiring
stock future gains nothing from moving the settlement price, because the gain on the final
margin is paid back in the price of the delivered shares. The incentive that motivated the
original design of this study, that holders of large expiring positions have a direct interest
in the settlement price, applies to stock futures only in accounting terms.

\subsection{Option strikes}

The incentive is real for stock options, because the settlement price determines which series
are exercised and therefore which option holders take or make delivery. If the settlement price
were drawn toward strikes on expiry days, either through the hedging of option writers or
through trading by option holders, the distance between the settlement price and the nearest
strike would be smaller on the expiry day than on other days of the month.

\begin{table}[htbp]
\centering
\small
\caption{Distance of the settlement price to the nearest option strike, in units of the local
strike interval. Other sessions are the sessions of the same expiry month; regressions include
a fixed effect for each security and month. High option activity is option volume on the
expiry and control days above the median relative to the security's cash turnover. The change
in the window is the distance at the settlement price less the distance at the pre-window
price. Standard errors clustered by session in parentheses.}
\label{tab:pinning}
\resizebox{\textwidth}{!}{\input{generated/tables/pinning}}
\end{table}

Table~\ref{tab:pinning} finds no such attraction. On other sessions of the month the average
distance is \PinDistanceControlmean{} strike intervals, close to the value of one quarter
expected when prices are unrelated to strikes. On expiry days the difference is
$\PinDistanceExpiry{}$ intervals (standard error \PinDistanceExpirySe{}), which is small and
not significant, and the share of settlement prices within a tenth of an interval of a strike
does not rise. Nor does the settlement window move the price toward the nearest strike: the
change in distance over the window is $\PinChangeExpiry{}$ on expiry days relative to other
sessions. For securities with the most option activity, where the incentive is greatest, the
estimates point in the direction opposite to attraction. Figure~\ref{fig:pinning} shows the
distribution.

\begin{figure}[htbp]
  \centering
  \includegraphics[width=0.62\textwidth]{pinning.pdf}
  \caption{Distribution of the distance of the settlement price to the nearest strike on expiry
    sessions and on the other sessions of the same month. The dashed line is the uniform
    density.}
  \label{fig:pinning}
\end{figure}

\subsection{Index strikes}

The index options have strikes at every fifty points, and their settlement is in cash. The
distance of the Nifty 50's settlement value to the nearest multiple of fifty, in units of the
interval, averages \IdxpinNiftyDistsettleNormal{} on sessions without an expiry, the value
expected when the index is unrelated to the strikes. On index expiries it differs by
$\IdxpinNiftyDistsettle{}$ (standard error \IdxpinNiftyDistsettleSe{}), and the change in
distance over the window by $\IdxpinNiftyDistchange{}$. The downward displacement of the index
settlement documented in Section~\ref{sec:index} is therefore not a movement toward the
nearest strike. It is a movement in one direction.
